\documentclass[11pt]{article}
\usepackage{geometry}\geometry{margin=1in}
\usepackage{hyperref}
\begin{document}

\begin{flushright}
\today
\end{flushright}

\vspace{1cm}

\noindent
Dear Editor,

\vspace{0.5cm}

\noindent
We submit the manuscript entitled ``\textbf{Cooperative Yaw Control for Wind Farm Power Maximization Using a Calibrated Gray-Box Wake Model and Reinforcement Learning}'' for consideration for publication in \emph{Wind Energy}.

\vspace{0.3cm}

\noindent
\textbf{Relevance to Wind Energy.}
The manuscript addresses a core wind-farm control problem---cooperative yaw steering for wake-induced power loss mitigation---using a control-oriented engineering methodology. The contribution lies at the intersection of wake modeling (calibrated Bastankhah--Port\'e-Agel Gaussian model, jointly fitted on three CFD datasets and cross-validated against FLORIS and Horns Rev~I LES), closed-loop yaw control (real-time policy via reinforcement learning with sub-millisecond inference), and actuator-aware dynamic operation (gated regime-conditional activation with yaw deadband). We believe this aligns with \emph{Wind Energy}'s readership of wind-farm operators, control engineers, and wake-modeling researchers.

\vspace{0.3cm}

\noindent
\textbf{Problem and approach.}
Wake-induced power losses motivate cooperative yaw control, yet conventional static optimization produces yaw vectors tied to single inflow snapshots and must be re-solved offline---a pipeline incompatible with the second-to-minute variability of the atmospheric boundary layer. We train a closed-loop reinforcement-learning policy on a CFD-calibrated gray-box wake model. The policy produces actuator-feasible yaw commands at sub-millisecond latency and is evaluated under steady-state, dynamic-wind, and noise-perturbed conditions.

\vspace{0.3cm}

\noindent
\textbf{Key findings.}
\begin{itemize}
\item The calibrated wake model reduces two-turbine yaw-case power-prediction error from $27.6\%$ to $6.6\%$, validated against Horns Rev~I LES.
\item Under steady-state evaluation the learned policy achieves $+4.91\%$ farm-power gain in wake-aligned regimes, recovering $94.9\%$ of the SLSQP static optimum (bootstrap 95\% CI $[93.2\%, 95.9\%]$).
\item FLORIS cross-validation on $3{,}000$ conditions with five independent seeds yields $+5.02\%$ aligned-cube gain (4.1\% erosion relative to gray-box), confirming conservative gain estimates.
\item Under dynamic AR(1) wind, a gated activation strategy with $2^\circ$ industrial yaw deadband reduces cumulative yaw travel by $43\%$ while retaining cooperative gains.
\item Observation-noise tests (wind direction up to $10^\circ$, wind speed up to $1.0$\,m/s, yaw angle up to $2^\circ$) show negligible performance degradation.
\item A Behavioral Cloning baseline that directly imitates the SLSQP oracle fails catastrophically ($-20.20\%$ marginal gain), confirming that the learned policy is structurally distinct from supervised imitation.
\end{itemize}

\vspace{0.3cm}

\noindent
\textbf{Scope and limitations.}
This is a simulation and cross-validation study. The gray-box model is control-oriented; FLORIS serves as an independent benchmark. No field SCADA validation or aeroelastic load assessment has been performed. These limitations are stated explicitly in the manuscript.

\vspace{0.3cm}

\noindent
\textbf{Originality.}
This work is original, has not been published elsewhere, and is not under consideration by another journal. All authors have approved the manuscript and this submission.

\vspace{0.5cm}

\noindent
\textbf{Corresponding Author:}\\
Zhuang Shao\\
China Resources Power Technology Research Institute Co., Ltd.\\
Shenzhen 518000, Guangdong Province, China\\
E-mail: shaozhuang@crpower.com.cn\\
ORCID: 0000-0003-2496-0797

\vspace{0.3cm}

\noindent
\textbf{Co-authors:}\\
Lijun Lei (Rundian Energy Science and Technology Co., Ltd.)\\
Peng Wang (North China University of Water Resources and Electric Power)\\
Liang Zheng (Rundian Energy Science and Technology Co., Ltd.)\\
Jie Zhou (Rundian Energy Science and Technology Co., Ltd.)

\vspace{0.3cm}

\noindent
\textbf{Suggested Reviewers:}
\begin{itemize}
\item Jianxin Zhou (zjx@seu.edu.cn), Southeast University
\item Zhenlong Wu (wuzhenlong2020@zzu.edu.cn), Zhengzhou University
\item Cong Yu (congy@jhun.edu.cn), Jianghan University
\item Hui Gu (guhuini@126.com)
\end{itemize}

\vspace{0.5cm}

\noindent
Sincerely,\\[0.3cm]
Zhuang Shao\\
(on behalf of all authors)

\end{document}
